
Spatial molecular profiling measures molecular state and position in the same tissue, and many methods now use position to cluster cells or capture spots into domains. A recurring question is whether a learned spatial encoder captures tissue organization that a simpler use of position would miss. We ask this question for one released graph encoder, Tessera, by asking how much of its agreement with supplied tissue annotations is explained by averaging each observation with its spatial neighbors at an appropriate scale.

The answer depends on what the annotation measures. Cell identity, neighborhood composition and tissue compartment are related but distinct targets: a molecular cell type can occur in several compartments, and a compartment can contain many cell types \citep{moses2022museum,rao2021architecture}. Spatial smoothing can therefore help recover a compartment while obscuring the cell types inside it. Tumor microenvironments make the distinction concrete because immune, stromal and epithelial populations intermix. Imaging-based panels, registered multi-assay breast measurements and serial-section Open-ST data each support several annotation types \citep{he2022cosmx,janesick2023breast,schott2024openst}. A region computed from neighboring cell types, a molecular category transferred from a reference and a boundary drawn by a pathologist answer different questions, so the origin of the annotation is part of the evaluation.

Spatial methods differ in graph construction, objective and assignment rule. STAGATE uses an adaptive graph-attention autoencoder, SEDR combines expression and spatial representations, GraphST applies graph self-supervised contrastive learning, and SpaGCN and SpaceFlow encode spatial consistency in other ways \citep{dong2022stagate,xu2024sedr,long2023graphst,hu2021spagcn,ren2022spaceflow}. Simpler approaches are competitive. SpatialPCA builds spatial correlation into dimension reduction, BANKSY augments each cell's features with a weighted neighborhood average, and MENDER summarizes multi-scale cellular context \citep{shang2022spatialpca,singhal2024banksy,yuan2024mender}. BANKSY in particular makes the balance between a cell and its neighborhood an explicit, tunable choice that spans cell typing and domain segmentation. Benchmarks compare many methods across technologies and tissues \citep{yuan2024benchmark,kang2025nar,chen2025imeta}, but a comparison of final partitions across methods conflates input processing, graph construction and clustering.

To attribute a gain to an encoder, the molecular input and assignment rule must be held fixed while the use of position is varied. Two further controls are needed. First, a simple spatial baseline must be given the same opportunity to choose its scale as the learned model has to fit its parameters; an untuned six-neighbor average is not that baseline. Second, the reference annotation's own construction must be considered. The lung regions used here were computed from cell-type frequencies among 200 spatial neighbors \citep{he2022cosmx,sun2026consensus}, so averaging at a similar scale is expected to recover them without any learned component.

The author-developed Tessera implementation evaluated here uses the released neural core with the field-aware graph and data adapters described below. Tessera combines learned edge gates, residual message passing, attention over network depths, feature reconstruction and a gate-weighted contrastive objective. Its gates were designed to down-weight edges that cross tissue boundaries. Each component makes a testable prediction: if the gates detect boundaries, they should separate annotation seams better than a trivial edge score and their removal should reduce agreement; if message passing captures structure beyond smoothing, Tessera should exceed a tuned neighbor average; and if the representation generalizes, it should help in patients that did not fit it. Coordinate identity and evaluation units also matter. Imaging fields can have separate local coordinate systems \citep{palla2022squidpy,yan2026alignment}, and thousands of cells from one field do not replace independent patients \citep{zimmerman2021pseudoreplication,roberts2017crossvalidation}.

We evaluated 13 assay objects: five NSCLC patients, four BCC patients, registered breast Visium, colorectal tubulovillous adenoma, an Open-ST lymph-node section and MIBI fields. Nine fixed arms, 29 spatial scale-control configurations per object plus PCA50, whole-section reruns, a corrected geometry null and patient-held-out transfer were scored with the same molecular input and assignment rule (Fig.~\ref{fig:workflow}). We report every arm and seed, use patients as the unit of inference, and test each component prediction directly.

Two further sets of data test whether these conclusions travel. An independent cohort of eight HER2-positive breast tumors carries regions drawn by a pathologist on the stained tissue \citep{andersson2021her2st}, a construction that involves no averaging of neighboring expression, and three non-tumor tissues with anatomical labels \citep{maynard2021dlpfc,wang2018starmap,moffitt2018merfish} show whether the ordering depends on tissue. The independent-cohort all-label target and the comparison with patient-excluded tuned controls were fixed in writing before fitting. The untuned same-graph comparator, weighting controls, boundary variants and tissue panel were added after the initial results and are exploratory.



